\section{PSM design regions: flux, saliency, and characteristic current}
\label{sec:psm-design-map}

\subsection{Nonsalient machine: two closed design inequalities}

For \(L_d=L_q\), positive torque is
\(M=k_T\psi_{\mathrm{pm}}i_q\).  On the loss disk
\(a(i_d^2+i_q^2)\leq P_{\mathrm{Cu},\max}\), the largest \(i_q\) is
\(I_P=\sqrt{P_{\mathrm{Cu},\max}/a}\), attained at \(i_d=0\).  Hence
\begin{equation}
 \boxed{M_{\max}^{(U\,\mathrm{free})}
 =k_T\psi_{\mathrm{pm}}\sqrt{\frac{P_{\mathrm{Cu},\max}}{a}}.}
 \label{eq:psm-nonsalient-mmax}
\end{equation}
For \(\psi_{\mathrm{pm}}>0\), the exact loss-limited inverse is
\begin{equation}
 \boxed{\psi_{\mathrm{pm}}\geq\psi_{\min}
 :=\frac{M_{\mathrm{req}}}{k_T}
 \sqrt{\frac{a}{P_{\mathrm{Cu},\max}}}.}
 \label{eq:psm-minimum-flux}
\end{equation}
This is a necessary and sufficient low-speed condition only while voltage is
inactive and the linear, nonsalient model is valid.

The characteristic current is \(I_{\mathrm{ch}}=\psi_{\mathrm{pm}}/L_d\).
Requiring the zero-flux point \((-I_{\mathrm{ch}},0)\) to lie inside the loss
circle gives
\begin{equation}
 \boxed{\psi_{\mathrm{pm}}\leq\psi_{\max}^{\mathrm{FW}}
 :=L_d\sqrt{\frac{P_{\mathrm{Cu},\max}}{a}}.}
 \label{eq:psm-characteristic-design}
\end{equation}
Combining the two yields the analytic window
\begin{equation}
 \boxed{\frac{M_{\mathrm{req}}}{k_T}
 \sqrt{\frac{a}{P_{\mathrm{Cu},\max}}}
 \leq\psi_{\mathrm{pm}}\leq
 L_d\sqrt{\frac{P_{\mathrm{Cu},\max}}{a}}.}
 \label{eq:psm-design-window}
\end{equation}
It is nonempty exactly when
\begin{equation}
 M_{\mathrm{req}}\leq\frac{k_TL_d}{a}P_{\mathrm{Cu},\max}.
 \label{eq:psm-window-nonempty}
\end{equation}
The upper bound is not a universal necessary condition for every finite-speed
requirement.  It is the familiar ideal flux-cancellation criterion associated
with wide or theoretically unbounded constant-power operation when resistance,
saturation, demagnetization, and other constraints are neglected
\cite{soong1994,morimoto1990}.  Modern design work still treats characteristic
current as a field-weakening metric, but evaluates it together with losses and
the full operating envelope \cite{lee2017,li2026}.

\begin{figure}[tb]
\centering
\begin{tikzpicture}
\begin{axis}[
 width=.78\linewidth,height=.50\linewidth,
 xmin=0,xmax=4.8,ymin=0,ymax=4.2,
 xlabel={$L_d$},ylabel={$\psi_{\mathrm{pm}}$},
 axis lines=left,ticks=none,clip=false]
 \addplot[name path=lower,red,very thick,domain=0:4.8] {1.25};
 \addplot[name path=upper,blue,very thick,domain=0:4.8] {.8*x};
 \addplot[name path=top,draw=none,domain=0:4.8] {4.2};
 \addplot[green!35] fill between[of=lower and upper,soft clip={domain=1.5625:4.8}];
 \node[red,anchor=west,font=\scriptsize] at (axis cs:3.35,1.05)
   {$\psi_{\min}$: torque};
 \node[blue,anchor=west,font=\scriptsize,rotate=32] at (axis cs:3.05,2.65)
   {$\psi=L_dI_P$: characteristic current};
 \node[green!40!black,font=\small] at (axis cs:3.1,1.85) {feasible window};
 \node[red!70!black,font=\scriptsize] at (axis cs:1,.65) {insufficient torque};
 \node[blue!70!black,font=\scriptsize,align=center] at (axis cs:2.2,3.45)
   {zero-flux point\\outside loss circle};
 \addplot[only marks,mark=*,mark size=2.7pt,black] coordinates {(2.7,1.7)};
\end{axis}
\end{tikzpicture}
\caption{Analytic nonsalient PSM design window.  The minimum PM flux follows
from low-speed torque; the upper ray places characteristic current inside the
copper-loss circle.  The latter is an ideal field-weakening criterion, not a
complete finite-speed feasibility proof.}
\label{fig:psm-design-window}
\end{figure}


Too little PM flux fails low-speed torque.  Increasing flux improves torque per
ampere, until the fixed flux consumes too much of the current budget when it
must be cancelled at high speed.  In coordinates
\((L_d,\psi_{\mathrm{pm}})\), the torque condition is a horizontal line and
the characteristic-current condition is a ray.  Their intersection makes the
tradeoff visible without solving a current trajectory.

\subsection{Salient PSM: a two-descriptor extension}

For a salient machine and \(\psi_{\mathrm{pm}}>0\), write a loss-boundary
current as \(i_d=I_P\cos\gamma\), \(i_q=I_P\sin\gamma\).  Then
\begin{equation}
 \frac{M}{k_T\psi_{\mathrm{pm}}I_P}
 =\sin\gamma\left(1+\gls{sym:sigma}\cos\gamma\right),
 \qquad \sigma=\frac{\Delta L I_P}{\psi_{\mathrm{pm}}}.
 \label{eq:psm-normalized-torque}
\end{equation}
Stationarity gives
\begin{equation}
 \cos\gamma+\sigma(2\cos^2\gamma-1)=0.
\end{equation}
For \(\sigma\neq0\), the physical maximum continuously connected to
\(\gamma=\pi/2\) has
\begin{equation}
 c_\star:=\cos\gamma_\star
 =\frac{-1+\sqrt{1+8\sigma^2}}{4\sigma},
 \quad
 F(\sigma)=\sqrt{1-c_\star^2}(1+\sigma c_\star),
 \label{eq:salient-factor}
\end{equation}
and \(F(0)=1\) by continuity.  Therefore
\begin{equation}
 M_{\max}^{(U\,\mathrm{free})}
 =k_T\psi_{\mathrm{pm}}I_PF(\sigma).
 \label{eq:salient-psm-capability}
\end{equation}
This is an analytic scalar design test, although inversion for
\(\psi_{\mathrm{pm}}\) is implicit because \(\sigma\) contains that flux.

A second independent coordinate is
\begin{equation}
 \gls{sym:iota}=\frac{I_{\mathrm{ch}}}{I_P}
 =\frac{\psi_{\mathrm{pm}}}{L_dI_P}.
 \label{eq:characteristic-ratio}
\end{equation}
The loss-limited torque shape depends on \(\sigma\), whereas ideal
flux-cancellation access is the half-plane \(\iota\leq1\).  Since
\(\sigma=(1-L_q/L_d)/\iota\), the saliency ratio and normalized PM flux are
not an arbitrary list of indices: they are alternative coordinates generated
by the torque circle and characteristic-current line.  Machines with the same
maximum torque per loss can still have different \(\iota\), voltage demand,
and high-speed behavior.  One capability number is therefore insufficient.

